% ID: SEL-D0024-OA-E01
% ESTADO: cotejada. Segunda pasada visual completa y borrador compilado revisados.
% Fuente: Selectividad/matematicas_aplicadas/reserva_b_EXAMEN 2-ANDALUCÍA.pdf
% SHA-256: 98b7723232acd3143de28397f7641923ed48292db7c35883ebd2b56a72b30b79
% Página: 1; opción A; ejercicio original 1.

\textbf{(2.5 puntos)} Un fabricante de complementos alimenticios elabora dos tipos de bebidas energéticas a partir de tres componentes: taurina, cafeína y L-carnitina. Un envase del primer tipo de bebida precisa 30 g de taurina, 40 g de cafeína y 20 g de L-carnitina, mientras que uno del segundo necesita 40 g de taurina, 30 g de cafeína y 10 g de L-carnitina. Sabiendo que dispone de 52 kg de taurina, 46 kg de cafeína y 20 kg de L-carnitina, que cada envase del primer tipo se vende por 1.5 \texteuro{} y cada envase del segundo tipo por 1 \texteuro{}, \textquestiondown{}cuántos envases de cada tipo de bebida tendría que elaborar para obtener la ganancia máxima? \textquestiondown{}A cuánto ascendería esta ganancia?
